\documentclass{article}
\usepackage{amsmath}
\usepackage[margin=1in]{geometry}
\usepackage{amssymb}

\begin{document}
\newpage
\section{Problem \#7}

Consider a counterflow, concentric tube heat exchanger that is designed to heat water from 20 to 80 \( ^{\circ}\)C. Hot oil is used as the heating agent, which enters the annulus at 160 \( ^{\circ}\)C and leaves at 140 \( ^{\circ}\)C. The inner tube, through which water flows, is thin-walled with a diameter of 20 mm. The overall heat transfer coefficient is given as 500 W/m\(^2\)K, and the design calls for a total heat transfer rate of 3000 W.

\subsection*{a. Determining the Length of the Heat Exchanger}

To find the required length of the heat exchanger, we first need to calculate the logarithmic mean temperature difference (LMTD). The LMTD is a correction factor that accounts for the temperature variation across the heat exchanger and is defined as:

\begin{equation}
\Delta T_{lm} = \frac{\Delta T_{o} - \Delta T_{i}}{\ln\left(\frac{\Delta T_{o}}{\Delta T_{i}}\right)}
\end{equation}

where:
\begin{itemize}
    \item \( \Delta T_{o} \) is the temperature difference between the hot and cold fluids at the outlet.
    \item \( \Delta T_{i} \) is the temperature difference at the inlet.
\end{itemize}

Given:
\begin{itemize}
    \item Inlet temperature of hot oil, \( T_{h,i} = 160 \, ^{\circ}C \).
    \item Outlet temperature of hot oil, \( T_{h,o} = 140 \, ^{\circ}C \).
    \item Outlet temperature of water, \( T_{c,o} = 80 \, ^{\circ}C \).
    \item Inlet temperature of water, \( T_{c,i} = 20 \, ^{\circ}C \).
\end{itemize}

We can now calculate the temperature differences:
\begin{equation}
\Delta T_{i} = T_{h,i} - T_{c,o} = 160 - 80 = 80 \, ^{\circ}C
\end{equation}
\begin{equation}
\Delta T_{o} = T_{h,o} - T_{c,i} = 140 - 20 = 120 \, ^{\circ}C
\end{equation}

Substituting into the LMTD equation, we get:
\begin{equation}
\Delta T_{lm} = \frac{120 - 80}{\ln(120/80)} = 98.65 \, ^{\circ}C
\end{equation}

Using the heat transfer rate equation \( q = U \times A \times \Delta T_{lm} \), and solving for the area \( A \), we can find the length \( L \) since the area is also a function of length (\( A = \pi D_i L \)):
\begin{equation}
L = \frac{q}{U \times \pi \times D_{i} \times \Delta T_{lm}} = \frac{3000 \, \text{W}}{500 \, \text{W/m}^2 \text{K} \times \pi \times 0.02 \, \text{m} \times 98.65 \, \text{K}} = 0.968 \, \text{m}
\end{equation}

\subsection*{b. Performance Degradation Due to Fouling}

After three years, fouling has occurred on the water side, reducing the outlet temperature of the water to 65 \( ^{\circ}\)C. Assuming the flow rates and inlet temperatures of both fluids remain unchanged, we now have to calculate the new heat transfer rate \( q_{new} \), the outlet temperature of the oil \( T_{h,o,new} \), the new overall heat transfer coefficient \( U_{new} \), and the fouling factor \( R''_{f,c} \).

First, we update the outlet temperature of the water in the heat transfer rate equation:
\begin{equation}
q_{new} = \dot{m}_{c} C_{p,c} (T_{c,i,new} - T_{c,i}) = 50 \times (65 - 20) = 2250 \, \text{W}
\end{equation}

Next, we solve for the new outlet temperature of the oil:
\begin{equation}
T_{h,o,new} = T_{h,i} - \frac{q_{new}}{\dot{m}_{h} C_{p,h}} = 160 - \frac{2250 \, \text{W}}{150 \, \text{W/K}} = 145 \, ^{\circ}C
\end{equation}

We then recalculate the temperature differences and the LMTD with the new temperatures:
\begin{equation}
\Delta T_{i,new} = T_{h,i} - T_{c,o,new} = 160 - 65 = 95 \, ^{\circ}C
\end{equation}
\begin{equation}
\Delta T_{o,new} = T_{h,o,new} - T_{c,i} = 145 - 20 = 125 \, ^{\circ}C
\end{equation}
\begin{equation}
\Delta T_{lm,new} = \frac{\Delta T_{o,new} - \Delta T_{i,new}}{\ln(\Delta T_{o,new}/\Delta T_{i,new})} = \frac{125 - 95}{\ln(125/95)} = 109.31 \, ^{\circ}C
\end{equation}

Using the new LMTD, we can solve for the new overall heat transfer coefficient:
\begin{equation}
U_{new} = \frac{q_{new}}{A \times \Delta T_{lm,new}} = \frac{2250 \, \text{W}}{\pi \times 0.02 \, \text{m} \times 0.968 \, \text{m} \times 109.31 \, \text{K}} = 338.42 \, \text{W/m}^2\text{K}
\end{equation}

Finally, we can find the water-side fouling factor by rearranging the equation for the overall heat transfer coefficient and solving for \( R''_{f,c} \):
\begin{equation}
R''_{f,c} = \frac{1}{U_{new}} - \frac{1}{U} = \frac{1}{338.42 \, \text{W/m}^2\text{K}} - \frac{1}{500 \, \text{W/m}^2\text{K}} = 0.00095489 \, \text{m}^2\text{K/W}
\end{equation}

This represents the additional thermal resistance due to fouling. It is a measure of the decrease in performance of the heat exchanger due to the buildup of material on the water side that inhibits heat transfer.
\end{document}
